\documentclass[11pt]{article}
\usepackage[margin=25mm]{geometry}
\usepackage{fontspec}
\setmainfont{Latin Modern Roman}
\usepackage{amsmath,amssymb,mathtools}
\usepackage{unicode-math}
\setmathfont{Latin Modern Math}
\usepackage{xurl,hyperref}
\hypersetup{colorlinks=true,urlcolor=blue}
\setcounter{secnumdepth}{0}
\setlength{\parindent}{0pt}
\setlength{\parskip}{5pt}
\setlength{\emergencystretch}{3em}
\newcommand{\sref}[2]{\hyperlink{source:#1:#2}{[S#2]}}
\newcommand{\cataloguescope}[1]{\paragraph{Recorded target.}#1}
\begin{document}
\section{Hadwiger Conjecture (Covering)}\label{statement}

\subsection{Definitions and mathematical statement}

\paragraph{Shared notation.}$\mathbb N=\{1,2,\ldots\}$, and $\mathbb Z,\mathbb Q,\mathbb R,\mathbb C$ denote the integers, rationals, reals and complex numbers. Any other conventions are specified locally.


A convex body $K\subset\mathbb R^n$ is compact and convex with nonempty interior. A smaller positive homothetic copy is $a+\lambda K=\{a+\lambda x:x\in K\}$, where $a\in\mathbb R^n$ and $0<\lambda<1$. Copies may have different translation vectors and contraction factors; rotations and negative homotheties are not allowed.
\paragraph{Statement recorded in the supplied source.}
For every $n\in\mathbb N$ and every convex body $K\subset\mathbb R^n$, do there exist $m\le2^n$, translations $a_1,\ldots,a_m$ and factors $0<\lambda_j<1$ such that
\[
 K\subseteq\bigcup_{j=1}^{m}(a_j+\lambda_jK)?
\]
\sref{statement}{2}
\subsection{Short English statement}
Can every n-dimensional convex body be covered by at most 2 to the n translated, positively contracted copies of itself?
\subsection{Sources}
\begin{description}
\item[\hypertarget{source:statement:1}{S1}] ULAM.AI and the UnsolvedMath Contributors, \emph{UnsolvedMath: A Curated Collection of Open Mathematics Problems}, catalogue entry. CC BY 4.0. \url{https://huggingface.co/datasets/ulamai/UnsolvedMath}.
\item[\hypertarget{source:statement:2}{S2}] Yegor Gorodzha, \emph{Fractional illumination and the optimal exponential rate in Hadwiger's covering conjecture} (2026), arXiv:2609.23913v1, Section 1, first paragraph. \url{https://arxiv.org/abs/2609.23913v1}.
\end{description}
\label{end:statement}
\end{document}
